\section{Preliminaries and the unimodular reduction}
\label{pre:section}

Throughout the proof, all operator norms refer to counting-measure
$\ell^2(V)$. A kernel is called invariant if its value at $(gx,gy)$ equals
its value at $(x,y)$ for every automorphism $g$. For a vertex set $D$, the
same letter denotes the induced subgraph on $D$. Write $C_x^D$ for the
open cluster of $x$ in $D$, and put
\[
 C_x=C_x^V,\qquad s_x=|C_x|,\qquad
 T_p^D(x,y)=\Prob_p(x\leftrightarrow y\text{ in }D),\qquad T_p=T_p^V.
\]
Quantities requiring a vertex outside their domain are defined to be zero.
We use $\partial_E A$ for the set of edges with exactly one endpoint in
$A$, and $\Nb[A]$ for $A$ together with all its neighbors.

\subsection{Graph reductions and the two established inputs}
\label{pre:reductions}

\begin{lemma}[Reduction to simple graphs]\label{pre:simple-reduction}
It suffices to prove the operator gap and the infinite-cluster conclusion
for simple graphs satisfying the hypotheses of \cref{thm:main}.
\end{lemma}
\begin{proof}
Loops never affect connectivity and may be removed. Replace each remaining
edge of $G$, including each copy of a multiple edge, by a two-edge path
through its own new vertex, obtaining a simple graph $\widetilde G$.
Quasi-transitivity and local finiteness give a finite degree bound $d_0$
and finitely many edge orbits in $G$. The lifted automorphisms consequently
have finitely many vertex orbits in $\widetilde G$.

To verify nonamenability, let a finite set $A$ in $\widetilde G$ contain
$n$ old vertices and $m$ new vertices. Each distinct boundary neighbor in
$G$ of its old-vertex set supplies a distinct boundary vertex of $A$:
along a selected subdivided edge, use the new vertex if it is outside
$A$, and otherwise use the old boundary neighbor. Hence
$|\partial_V A|\ge h_V(G)n$. At most $d_0n$ selected new vertices have
an endpoint among the selected old vertices. If $m>2d_0n$, more than
$m/2$ selected new vertices have both endpoints outside $A$; bounded
degree then gives $|\partial_V A|\ge m/(2d_0)$. The first estimate when
$m\le2d_0n$ and the second otherwise give a positive lower bound for
$|\partial_V A|/(n+m)$, independent of $A$.

In percolation of parameter $t$ on $\widetilde G$, declare an original
edge open precisely when both its replacement edges are open. These
independent declarations have parameter $t^2$. Infinite clusters
correspond bijectively: a cluster with only finitely many old vertices
is finite by local finiteness. In particular
\[
 \pc(\widetilde G)^2=\pc(G),\qquad
 \pu(\widetilde G)^2=\pu(G).
\]
Moreover, $T_{t^2}^G$ is the old-vertex compression of
$T_t^{\widetilde G}$. Boundedness of the latter implies boundedness of
the former. Thus a strict operator gap, and a coupled interval of
infinitely many infinite clusters, both transfer to $G$.
\end{proof}

We henceforth assume that $G$ is simple and fix an integer $d$ bounding
its degrees. Write $\Adj(x,y)=\one_{\{x\sim y\}}$ for its adjacency kernel.
Positive vertex isoperimetry supplies a constant $\iota>0$
such that $|\partial_E A|\ge\iota|A|$ for every finite nonempty $A$.

\begin{lemma}[Critical finiteness]\label{pre:critical-finiteness}
We have $0<\pc<1$. At $\pc$, all open clusters are finite almost surely,
simultaneously at every vertex and in every induced subgraph under
restriction of the full configuration.
\end{lemma}
\begin{proof}
There are at most $d^n$ length-$n$ paths from a fixed vertex, so no
infinite cluster exists when $pd<1$. For the opposite bound, the number
of connected vertex sets of size $n$ containing a fixed vertex is at
most $d^{2n}$: choose a deterministic spanning tree of each set and
encode it by a traversal of length $2n-2$. If such a set is the root
cluster, all its boundary edges are closed. Consequently
\[
 \Prob_p(|C_x|<\infty)
 \le \sum_{n\ge1}d^{2n}(1-p)^{\iota n}<1
\]
when $p$ is sufficiently close to $1$.

For every graph-distance ball,
\[
 |B(x,r+1)|\ge(1+h_V(G))|B(x,r)|,
\]
so $G$ has exponential growth. Hutchcroft's critical-finiteness theorem
for connected locally finite quasi-transitive graphs of exponential
growth \cite[Theorem~1.2]{Hutchcroft2016} therefore applies to bond
percolation on $G$. Countability gives simultaneous finiteness at all
vertices. Every cluster in a restricted configuration is contained in
a full-graph cluster, proving the last assertion.
\end{proof}

\subsection{Product inequalities and operator bounds}
\label{pre:product-tools}

The first two assertions below are Harris's association inequality
\cite{Harris1960} and the van den Berg--Kesten (BK) inequality
\cite{VanDenBergKesten1985}. We include the product-measure proofs.

\begin{lemma}[Association, disjoint witnesses, and the tree bound]
\label{pre:probability-tools}
Increasing nonnegative functions of independent percolation bits have
nonnegative covariance. If increasing events have disjoint finite open
witnesses, their joint disjoint-occurrence probability is at most the
product of their probabilities. In particular, for all $x,y,z$,
\begin{equation}\label{eq:s01}
 \Prob_p(x\leftrightarrow y,\ x\leftrightarrow z)
 \le\sum_w T_p(x,w)T_p(w,y)T_p(w,z).
\end{equation}
The same assertions hold in every induced domain.
\end{lemma}
\begin{proof}
For association in a finite product, condition on the last bit and use
induction. The two conditional means are increasing functions of that
bit, whose covariance is nonnegative. Truncation and conditional
expectations on increasing finite coordinate sets extend this argument
to the countable product and to nonnegative functions.

Here is a finite-product proof of the disjoint-witness inequality.
For two increasing events $A,B$, denote disjoint occurrence by
$A\mathbin{\Box}B$. Suppose first that $p\le1/2$, and let $A_0,A_1$
and $B_0,B_1$ be the sections at the last bit. Write
$D_{ij}=A_i\mathbin{\Box}B_j$. The zero-section of
$A\mathbin{\Box}B$ is contained in $D_{00}$; its one-section is
contained in $D_{10}\cup D_{01}$. Also
$D_{00}\subseteq D_{10}\cap D_{01}$. Thus
\[
 \Prob(A\mathbin{\Box}B)
 \le(1-2p)\Prob(D_{00})+p\Prob(D_{10})+p\Prob(D_{01}).
\]
Apply induction, putting $a_i=\Prob(A_i)$ and $b_i=\Prob(B_i)$.
The difference between $\Prob(A)\Prob(B)$ and the resulting upper
bound is $p^2(a_1-a_0)(b_1-b_0)\ge0$.
For $1/2<p<1$, realize each bit as the logical OR of $m$ independent
bits of parameter $q=1-(1-p)^{1/m}\le1/2$. Disjoint open witnesses
lift to disjoint open witnesses in the enlarged product. The cases
$p=0,1$ are immediate. Iteration proves the assertion for any finite
number of events, and exhaustion proves the finite-witness assertion
on $G$.

On $\{x\leftrightarrow y,\ x\leftrightarrow z\}$, take a finite open
tree joining the three vertices. Its branch vertex $w$ supplies three
edge-disjoint paths to $x,y,z$, with zero-length paths permitted.
A union bound over $w$, followed by the disjoint-witness inequality,
gives \eqref{eq:s01}.
\end{proof}

The tree bound is the three-terminal instance of the tree-diagram
method of Aizenman and Newman \cite{AizenmanNewman1984}.
We repeatedly use two elementary kernel facts. If $|A(x,y)|\le B(x,y)$
and $B$ is a nonnegative bounded kernel, then $\norm A\le\norm B$:
test the bilinear form against absolute values of finitely supported
vectors. If all row sums and column sums of $B$ are at most $b$, then
$\norm B\le b$, by Cauchy--Schwarz with the weights $B(x,y)$.

\begin{lemma}[The operator threshold precedes uniqueness]
\label{pre:threshold-order}
For every quasi-transitive graph under consideration,
$\ptwo\le\pu$.
\end{lemma}
\begin{proof}
Suppose that there is a unique infinite cluster almost surely at
parameter $p$. Some vertex has positive probability of belonging to
it, by countability. Association with the event that a fixed finite
connecting path is open makes this probability positive at every
vertex. There are finitely many vertex types, so
$\inf_x\Prob_p(|C_x|=\infty)=\theta>0$. Association and uniqueness give
$T_p(x,y)\ge\theta^2$ for every pair $x,y$. For finite $F\subset V$,
\[
 \ip{\one_F}{T_p\one_F}\ge\theta^2|F|^2,
\]
which rules out boundedness on $\ell^2(V)$. Boundedness parameters
form a downward interval by entrywise domination. Thus $\ptwo$ is
at most every parameter of almost-sure uniqueness, proving the claim
under the stated supremum and infimum definitions.
\end{proof}

\begin{proposition}[Reduction to a unimodular action]
\label{pre:unimodular-reduction}
If $\Aut(G)$ is nonunimodular, then $\pc<\ptwo\le\pu$ and
$T_\pc$ is bounded.
For the proof of the operator gap it therefore remains to treat the
case in which $\Gamma=\Aut(G)$ is unimodular.
\end{proposition}
\begin{proof}
Give $\Aut(G)$ the topology of pointwise convergence on vertices.
Vertex stabilizers are open and compact: restrictions to successive
finite balls give the usual diagonal compactness argument. Thus
$\Aut(G)$ is a locally compact group. If it is nonunimodular, its
quasi-transitive action satisfies the hypotheses of
\cite[Theorem~2.9]{Hutchcroft2019}, which gives
$\pc<p_{q\to q}$ for every $1<q<\infty$. Take $q=2$ and apply
\cref{pre:threshold-order}. A bounded kernel at any parameter strictly
between $\pc$ and $\ptwo$ dominates $T_\pc$, proving its boundedness.
This cited theorem requires a
quasi-transitive nonunimodular automorphism subgroup; hyperbolicity
is not among its hypotheses.
\end{proof}

The remaining assertion, proved in \cref{con:section}, is the principal
new estimate.

\begin{theorem}[Critical boundedness for a unimodular action]
\label{thm:critical-operator}
Let $G$ be an infinite connected locally finite simple quasi-transitive
graph with $h_V(G)>0$. If $\Aut(G)$ is unimodular, then
$\norm{T_{\pc}}_{2\to2}<\infty$.
\end{theorem}

\subsection{Averaging over types and a finite trace}
\label{pre:types}

The inverse stabilizer-volume weights below give the quasi-transitive
mass-transport principle of \cite[Corollary~3.5]{BLPS1999Invariant}; see
also \cite[Theorem~3.1]{AldousLyons2007} and
\cite[Corollary~8.11 and Section~8.2]{LyonsPeres2016}.

For the rest of the main argument assume that $\Gamma$ is unimodular.
Let $O_1,\ldots,O_{h_0}$ be its vertex orbits, with representatives
$o_1,\ldots,o_{h_0}$. Fix a Haar measure on $\Gamma$, and set
\[
 \sigma_i=
 \frac{\operatorname{vol}(\Gamma_{o_i})^{-1}}
 {\sum_{j=1}^{h_0}\operatorname{vol}(\Gamma_{o_j})^{-1}},
 \qquad [F(o)]=\sum_{i=1}^{h_0}\sigma_iF(o_i),
 \qquad c_\sigma=(\min_i\sigma_i)^{-1}.
\]
Brackets always mean this finite type average, including when $F$
already contains a probabilistic expectation. These weights determine a
law on the root types, not a measure replacing counting measure on $V$.
In particular all adjoints, positivity assertions, and operator norms
below are those of the previously defined Hilbert space $\ell^2(V)$.

\begin{lemma}[Weighted mass transport]\label{pre:mass-transport}
For every invariant nonnegative function $m$ on ordered vertex pairs,
\begin{equation}\label{eq:s02}
 \left[\sum_y m(o,y)\right]
 =\left[\sum_x m(x,o)\right].
\end{equation}
The identity also applies to the expectation of a covariant random
mass and to absolutely summable signed or complex masses.
In particular, an invariant stochastic kernel on a single orbit has
column sums one, whether or not it is symmetric.
\end{lemma}
\begin{proof}
Decompose the ordered pairs into $\Gamma$-orbits. For one such orbit,
represented by $(x,y)$ of types $i,j$, its number of pairs with first
coordinate $x$ is
\[
 \frac{\operatorname{vol}(\Gamma_x)}
 {\operatorname{vol}(\Gamma_x\cap\Gamma_y)},
\]
and its number with second coordinate $y$ is the same expression with
$\Gamma_y$ in the numerator. Unimodularity makes the volumes of
conjugate compact subgroups equal. Multiplication by $\sigma_i$ and
$\sigma_j$ therefore gives the same contribution. Summing proves
\eqref{eq:s02}; Tonelli's theorem handles nonnegative infinite sums.
Take expectations for random masses, and use absolute convergence for
signed or complex ones. Supporting the mass on one orbit cancels its
positive type weight, giving the final assertion.
\end{proof}

The diagonal average gives a normalized trace even though $V$ is
infinite. This is the invariant-operator trace discussed in
\cite[Section~5]{AldousLyons2007}; we verify its precise form here.

\begin{lemma}[The invariant trace]\label{pre:trace}
For a bounded invariant operator $S$, define
$\tr S=[S(o,o)]$. This is a positive normalized trace:
$\tr I=1$ and $\tr(AB)=\tr(BA)$ for bounded invariant $A,B$.
If $A$ is positive semidefinite, invariant, and bounded, and $B=B^*$
is invariant and bounded,
then
\[
 |\tr(BA)|\le\norm B\,\tr A,
 \qquad A(x,x)\le c_\sigma\tr A.
\]
\end{lemma}
\begin{proof}
Positivity and normalization follow from the diagonal definition.
The diagonal product expansion is absolutely summable, since
\[
 \sum_y|A(x,y)B(y,x)|
 \le\norm{A^*\one_{\{x\}}}\,
       \norm{B\one_{\{x\}}}\le\norm A\norm B.
\]
Apply \eqref{eq:s02} to these products to obtain cyclicity.
The positive square root $A^{1/2}$ is invariant, because $A$ commutes
with the unitary action of every automorphism and functional calculus
preserves that commutation. The inequalities
\[
 -\norm B A\le A^{1/2}BA^{1/2}\le\norm B A
\]
and cyclicity prove the trace estimate. Finally the nonnegative
diagonal is constant on each type, so its value on type $i$ is at
most $\sigma_i^{-1}\tr A$.
\end{proof}

\subsection{A symmetric walk with a spectral gap}
\label{pre:walk-subsection}

We use a symmetric lazy version of the walk so that counting measure
is reversible even when degrees differ between vertex types. The
spectral estimate below is the discrete Cheeger argument; see
\cite{Dodziuk1984} for its isoperimetric background. Define
\[
 P(x,y)=\frac1{2d}\quad(x\sim y),\qquad
 P(x,x)=1-\frac{\deg(x)}{2d},
\]
with all other entries zero.

\begin{lemma}[Walk and type chain]\label{pre:walk}
The kernel $P$ is positive semidefinite and
\begin{equation}\label{eq:s03}
 \rho:=\norm P<1.
\end{equation}
Its projection onto the vertex types is an irreducible lazy finite
Markov chain. There exist $c_*>0$ and $l_0<\infty$ such that
\[
 \sum_{y\in O_j}P^l(x,y)\ge c_*
 \quad\text{for every }x\in V,\ j\le h_0,\ l\ge l_0.
\]
\end{lemma}
\begin{proof}
The kernel $2P-I$ is symmetric and stochastic, hence a contraction;
therefore $P$ is positive semidefinite. For a finitely supported
nonnegative $f$, the level-set identity and edge isoperimetry give
\[
 \sum_{\{x,y\}\in E}|f(x)^2-f(y)^2|
 =\int_0^\infty|\partial_E\{f^2>t\}|\,dt
 \ge\iota\norm f^2.
\]
Cauchy--Schwarz and
$\sum_{\{x,y\}\in E}(f(x)+f(y))^2\le2d\norm f^2$ imply
\[
 \ip f{(I-P)f}
 =\frac1{2d}\sum_{\{x,y\}\in E}(f(x)-f(y))^2
 \ge\frac{\iota^2}{4d^2}\norm f^2.
\]
Since $P$ has nonnegative entries,
$|\ip f{Pf}|\le\ip{|f|}{P|f|}$, so the same norm bound extends to
arbitrary finitely supported vectors and then to $\ell^2(V)$.

The transition probabilities between types are well defined by
invariance. Connectivity gives irreducibility, and every holding
probability is at least $1/2$. By following a connecting type path
and padding it with holds, a sufficiently large common power of this
finite transition matrix has every entry positive. Its minimum entry
is a lower bound for every later power, because left multiplication
by a stochastic matrix takes convex combinations of its rows.
\end{proof}

\subsection{Subcritical susceptibility, differentiation, and sprinkling}
\label{pre:subcritical}

Set
\[
 \chi_x(p)=\E_p|C_x|,\qquad
 \chi(p)=[\chi_o(p)],\qquad
 \chi_{\max}(p)=\sup_x\chi_x(p).
\]
The parameter will be suppressed when unambiguous.

\begin{lemma}[Subcritical finiteness and type comparison]
\label{pre:sharpness}
For $p<\pc$, $\chi_{\max}(p)<\infty$ and
\[
 \sup_x\E_p s_x^2\le\chi_{\max}(p)^3,
 \qquad \chi_{\max}(p)\le c_\sigma\chi(p).
\]
There is $c>0$, depending only on $G$, such that
$\chi_x(p)\ge c\chi(p)$ for every $x$ and
$\pc/2\le p<\pc$. The operator $T_p$ is symmetric, positive
semidefinite, and bounded with $\norm{T_p}\le\chi_{\max}(p)$.
\end{lemma}
\begin{proof}
Susceptibility finiteness on quasi-transitive graphs is established in
\cite[Theorem~2]{AntunovicVeselic2008}, following the sharpness work of
Aizenman--Barsky and Menshikov \cite{AizenmanBarsky1987,Menshikov1986}.
We give a proof by adapting the finite-set
criterion of Duminil-Copin and Tassion
\cite[Section~1]{DuminilCopinTassion2016} to the finitely many vertex
types.
For a finite set $S\ni x$, put
\[
 \varphi_p(x,S)=p\sum_{\substack{a\sim b\\a\in S, b\notin S}}
                 \Prob_p(x\leftrightarrow a\text{ in }S),
\]
and let $p_*$ be the supremum of parameters for which every vertex has
such an $S$ with $\varphi_p(x,S)<1$. This set of parameters is
downward closed. If $p<p_*$, choose a qualifying larger parameter;
finitely many orbit representatives then supply translated test sets
with a common size bound $L$ and costs at $p$ bounded by $a<1$.

For a finite induced domain $\Lambda$, let $H_\Lambda$ be its largest
susceptibility. A simple open path from $x\in\Lambda$ to a vertex
outside its test set $S$ has a first exit edge from $S$. Its prefix
inside $S$, that edge, and its remaining leg have disjoint witnesses.
The disjoint-witness inequality gives
\[
 \E_p|C_x^\Lambda|\le L+\varphi_p(x,S)H_\Lambda
 \le L+aH_\Lambda.
\]
This remains valid if $S$ is not contained in $\Lambda$, since the
prefix inside $S\cap\Lambda$ is dominated by the connection in $S$.
Maximizing and exhausting $G$ gives $\chi_{\max}(p)\le L/(1-a)$.
Thus $\pc\ge p_*$.

For the reverse inequality, choose $p_*<r<p<1$. There is a vertex $x$
such that $\varphi_r(x,S)\ge1$ for every finite $S\ni x$. Let
$f_\Lambda(u)$ be the probability that $x$ connects outside a finite
set $\Lambda\ni x$. This is a finite-edge event. Finite-product
differentiation expresses its derivative as the expected number of
closed pivotal edges divided by $1-u$.
Explore from outside $\Lambda$, and let $S$ be the vertices of
$\Lambda$ not connected to its exterior. On disconnection $x\in S$.
Conditional on this exploration, the edges internal to $S$ remain
independent with parameter $u$, and its crossing edges are closed.
A crossing edge is pivotal precisely when its inner endpoint connects
to $x$ in $S$. Hence, for $r\le u\le p$,
\[
 f_\Lambda'(u)
 =\frac1{u(1-u)}
   \E_u\bigl[\one_{\{x\in S\}}\varphi_u(x,S)\bigr]
 \ge\frac{1-f_\Lambda(u)}{u(1-u)}.
\]
Integration gives
$f_\Lambda(p)\ge1-\exp(-\int_r^pdu/(u(1-u)))>0$, uniformly in
$\Lambda$. Exhaustion yields $\pc\le p$; letting $p\downarrow p_*$
proves $p_*=\pc$.

Sum \eqref{eq:s01} over $y,z$ to obtain the second moment bound.
The upper type comparison follows directly from the weights.
There is a common finite bound on the length needed to reach a vertex
of any prescribed type from any starting vertex: choose paths between
representatives and translate them. Association with such a path of
length $l$ gives $\chi_x(p)\ge p^l\chi_y(p)$, proving the lower
comparison on the stated parameter interval.
Symmetry of $T_p$ and the row/column sum bound prove boundedness.
For every finitely supported complex $f$,
\[
 \sum_{x,y}\overline{f(x)}T_p(x,y)f(y)
 =\E_p\sum_{C}\left|\sum_{x\in C}f(x)\right|^2\ge0,
\]
where the sum is over open clusters. This proves positive
semidefiniteness.
\end{proof}

\begin{lemma}[Russo identities in operator norm]\label{pre:russo}
The map $p\mapsto T_p$ is locally absolutely continuous in operator
norm on $(0,\pc)$, and $\chi$ is locally absolutely continuous there.
For almost every $p\in(0,\pc)$,
\begin{equation}\label{eq:s04}
 \begin{aligned}
 (1-p)T'_p(x,y)
 &=\sum_a\sum_{b\sim a}
   \Prob_p(x\in C_a,\ y\in C_b,\ C_a\ne C_b),\\
 (1-p)\chi'
 &=\left[\sum_{b\sim o}
   \E_p(s_o s_b;\ C_o\ne C_b)\right].
 \end{aligned}
\end{equation}
\end{lemma}
\begin{proof}
In a finite induced domain, differentiate the connection event by
closed pivotals. A closed pivotal edge has a unique orientation
$(a,b)$ for which $x\in C_a$ and $y\in C_b$; these two clusters are
distinct. This proves the finite-domain first identity with the
factor $1-p$ and no extra factor of two.

Fix a compact parameter interval contained in $(0,p_1]$ with
$p_1<\pc$. The summand is bounded, by disjoint witnesses, by
$T_{p_1}(x,a)T_{p_1}(b,y)$. Its sum over $a,b,y$ is at most
$d\chi_{\max}(p_1)^2$. Under exhaustion, connection events converge,
and distinctness indicators converge as well: if two vertices are
connected in the full graph, a finite path eventually witnesses this.
Dominated convergence therefore passes the integrated Russo identity
to the full graph.

For each representative $o_i$, its row is thus the integral of an
$\ell^1(V)$-valued integrand dominated coordinatewise by the fixed
summable row of $(1-p_1)^{-1}T_{p_1}\Adj T_{p_1}$. Lebesgue
differentiation gives row differentiation in $\ell^1$ almost
everywhere. One can see this directly by differentiating a finite
set of coordinates and bounding the remaining tails by the same
summable majorant. There are finitely many row types. Invariance,
symmetry, and the row/column sum bound consequently upgrade these
statements to local absolute continuity and differentiation in
operator norm. Summing the first identity over $y$ and then applying
\eqref{eq:s02} to move the root from $x$ to $a$ gives the second
identity. All sums are justified by the same domination.
\end{proof}

The next connection-operator comparison is
\cite[Lemma~2.4]{Hutchcroft2019}. We include the argument because it is
used both before and after the critical bound is proved.

\begin{lemma}[Sprinkling]\label{pre:sprinkling}
For $0\le p,a\le1$,
\begin{equation}\label{eq:s05}
 T_{p+(1-p)a}
 \le\sum_{j\ge0}a^jT_p(\Adj T_p)^j
 \qquad\text{entrywise}.
\end{equation}
If $T_p$ is bounded and $ad\norm{T_p}<1$, the kernel on the left is
bounded. In particular, boundedness of $T_\pc$ implies
$\pc<\ptwo$.
\end{lemma}
\begin{proof}
Take independent fields of parameters $p$ and $a$; their union has
parameter $p+(1-p)a$. A finite simple open path in the union can be
split at the distinct ordered edges for which it uses the second
field. The intervening first-field legs have disjoint open witnesses.
For prescribed legs and $j$ prescribed second-field edges, independence
and the disjoint-witness inequality give the corresponding product
of $j+1$ connection probabilities times $a^j$. Summing over all
intermediate vertices and $j$ yields \eqref{eq:s05}; extra paths and
repetitions in that sum only increase the bound. Since
$\norm{\Adj}\le d$, the right-hand series converges in operator norm
when $ad\norm{T_p}<1$. Entrywise domination gives the assertion.
Finally $\pc<1$ permits a positive $a$ satisfying this condition at
$p=\pc$.
\end{proof}
